% =============================================================================
% §1.7 重写草稿 v3（独立片段，未 patch 进主稿）
% 相对 v2 的结构性重写，针对两条专家级批评：
%   (A) 术语去「治理黑话」：节标题「实验设置与比较协议」→「实验设置」；
%       小节标题朴素化为领域通用语（训练目标 / 长期递推评估 / 评价指标 /
%       对比模型与结构消融 / 初值扰动与鲁棒性评估）；删除「协议轴 / 配置轴 /
%       协议敏感性 / 口径 / 协议单元」等来自内部溯源-审计文档的措辞。
%   (B) 详略翻正（按 DL/动力学论文的重要性排序）：
%       - 取消「三轴正交」命名框架与 tab:protocol-axes，改用领域散文表达
%         训练相/评估相、强度/采样方式的区分（红线「三者不混淆」仍达成）。
%       - 扩写「对比模型与结构消融」：把基线链条讲成隔离结构因素的仪器，
%         每个消融对应一条可检验结构假设并对接 §1.6 性质。
%       - 收紧「评价指标」的中位数+P95 理由（科学判断，值得正文）。
%       - 证据纳入/异常治理压成 1–2 句原则，机制下沉 §1.8（消与 §1.8 的重复）。
%       - 噪声数值留表、设计理由讲清；训练调度去 runbook 化。
% 红线：三轴不混淆；噪声为 filtered-state 鲁棒性正则、非导航后验；不报告性能数值/
%       排序/「最重要结构」结论；不宣称 v4lite 优越；§1.7 不预锁 §1.8 证据子集。
% 交叉引用沿用主稿既有 label；§1.8 仅引用 \label{sec:training-baselines-protocol}
%   （已保留），主稿现行 §1.7 各小节本无 label，故本节小节 label 可自由命名。
% =============================================================================

\section{实验设置}
\label{sec:training-baselines-protocol}

本节给出检验前述结构约束所用的实验设置。第~\ref{sec:structured-continuous-model} 节的参数化已在函数类层面保证正定逆质量、半正定耗散与斜对称零功率耦合，第~\ref{sec:energy-power-properties} 节给出了静水机械核心的功率平衡；这些性质由构造保证、不随训练改变。模型作为长期轨迹预测器的实际行为，则取决于数值积分、状态转换与数据分布的共同作用，须经训练、长期递推与诊断检验。本节依次说明训练目标（第~\ref{subsec:training-objective} 小节）、长期递推评估方式（第~\ref{subsec:rollout-eval} 小节）、评价指标（第~\ref{subsec:metrics} 小节）、用于隔离各结构因素的对比模型与消融设计（第~\ref{subsec:models-ablations} 小节），以及作为鲁棒性检验的初值扰动（第~\ref{subsec:ic-perturbation} 小节）。各结构因素的实际作用与定量结果在第~\ref{sec:results-structure-evidence} 节报告。

\subsection{训练目标}
\label{subsec:training-objective}

训练数据由控制块序列组成：在单个控制块内，控制命令与外源上下文（海流、深度）保持分段常值，模型从块初值积分生成预测轨迹，
\begin{equation}
  y_0=\mathcal{T}_{d\to m}(s_0),\qquad
  \hat y(t_i)=\Phi_\theta(t_i;t_0,y_0),\qquad
  \hat s_i=\mathcal{T}_{m\to d}(\hat y(t_i)),
  \label{eq:block-prediction}
\end{equation}
其中 \(\mathcal{T}_{d\to m}\) 与 \(\mathcal{T}_{m\to d}\) 为式~\eqref{eq:data-to-model-state} 与式~\eqref{eq:model-to-data-state} 的数据态--模型态转换，\(\Phi_\theta\) 为式~\eqref{eq:complete-augmented-vector-field} 向量场经神经常微分方程积分得到的流映射。目标轨迹同样转换到模型态，使有海流样本的速度监督作用于相对水速度 \(\nu_r\)（依式~\eqref{eq:velocity-contract} 的速度契约），位置、姿态与执行器状态在两种表示中保持同一数值。加权训练目标取
\begin{equation}
  \mathcal{L}
  =
  \sum_i
  \lambda_x\norm{\hat x_i-x_i}^2
  +\lambda_R d_{\SO(3)}(\hat R_i,R_i)^2
  +\lambda_\nu\norm{\hat\nu_{r,i}-\nu_{r,i}}^2
  +\lambda_a\norm{\hat u_{a,i}-u_{a,i}}^2
  +\lambda_{\SO}\mathcal{L}_{\SO(3)} .
  \label{eq:training-loss}
\end{equation}
姿态误差以旋转测地距离 \(d_{\SO(3)}(\hat R,R)\) 度量、而非旋转矩阵的欧氏差，使监督直接作用于 \(\SO(3)\) 流形上的几何偏离，并对脱离旋转群的非物理偏移保持敏感；\(\mathcal{L}_{\SO(3)}\) 为旋转矩阵的正交性与行列式正则项。速度项以相对水速度为监督对象，与第~\ref{sec:state-current-contract} 节的速度契约一致；评估时预测轨迹经 \(\mathcal{T}_{m\to d}\) 回到数据态，终端位置、姿态与总体速度等可观测量在数据空间解释（第~\ref{subsec:metrics} 小节）。连续时间积分、优化器与网络容量等实现细节统一列于表~\ref{tab:experimental-config}。

\subsection{长期递推评估}
\label{subsec:rollout-eval}

块级训练目标只反映局部受控初值问题的拟合质量，无法说明误差在自由递推中的累积。本文以长期自由递推为主要评估方式：将控制块依时间串接，首块由给定初值启动，其后每块以上一块的预测末状态为初值，并按数据协议更新位置原点、海流速度、控制命令与深度上下文。所得轨迹不再由真实状态校正，单块误差沿时间累积、放大，姿态漂移、速度发散与数值失败等仅在长期显现的问题由此暴露。评估使用受控仿真生成的 AUV 轨迹，区分无海流与有海流设置；有海流数据依式~\eqref{eq:velocity-contract} 维持总体速度与相对水速度的转换，使位姿几何按总体速度推进、水动力配对作用于相对水速度。

评估时域取 10、30 与 60~s 三层，以 60~s 为本章主时域；较短时域用于刻画误差随时间的增长形态与提前失败的位置，避免单一终端时刻掩盖中途已发散的轨迹。未完成目标时域的轨迹，其终端误差不计入该时域的条件误差统计，可比性由完成率反映。对未完成的轨迹，进一步区分模型预测发散与真值侧或数值侧的局部失败，使后续分析能把结构导致的长期不稳定与偶发的数值失败分开——这一区分对判断几何结构是否支撑长期稳定尤为关键。

\subsection{评价指标}
\label{subsec:metrics}

外部任务指标在数据空间定义，刻画模型作为长期轨迹预测器的可观测性能。主精度以 60~s 终端位置误差的中位数与第 95 百分位（P95）并列度量，分别刻画典型误差与尾部风险。采用并列双指标而非单一标量，是因为典型精度与尾部稳健并不总是同向：一个中位数很低的模型仍可能在少数轨迹上累积显著更大的误差，而单一中位数或单一均值都会掩盖这一点；其必要性将在第~\ref{sec:results-structure-evidence} 节由中位数与尾部排序不一致的具体情形体现。完成目标时域的轨迹另报告深度误差、旋转测地误差与总体线速度、角速度误差。完成率不参与精度排序，仅用于判断终端误差是否具可比性——终端误差只有在完成目标时域的样本上才有意义。

除外部指标外，本文报告若干模型空间的内部诊断量，用于说明各结构组件是否按设计工作，但不替代外部任务误差：相对水速度误差核验 \(\nu_r\) 是否与第~\ref{sec:state-current-contract} 节的速度契约一致；机械能量跨度与最大能量变化反映机械存储是否在递推中保持有限量级；\(\SO(3)\) 行列式与正交性误差反映姿态矩阵是否保持接近旋转群。各量先在单个运行内对评估轨迹取统计量，再在随机种子之间取均值，记号 \(N\) 表示参与聚合的有效种子数。

\subsection{对比模型与结构消融}
\label{subsec:models-ablations}

结构约束的作用通过一组受控对照检验。对比模型沿结构完整度递减排列，使每一步移除恰好对应一个结构因素：完整结构化模型 AUVHamNODE 同时具备 \(\SE(3)\) 位姿几何、机械存储与 \(D_\theta\)/\(J_\theta\)/\(\tau_\theta\) 分解；端口哈密顿风格基线保留机械存储核心，而改变非保守广义力的组织方式；\(\SE(3)\) 动量黑箱与加速度黑箱保留位姿几何乃至质量映射，却把动量或加速度动力学交由黑箱；全状态黑箱则不含任何显式几何或能量结构。由完整模型到全状态黑箱的链条因而把显式位姿几何、机械能量核心与非保守力分解等结构因素逐层剥离，使各因素的作用可在对照中分离。其中 \(\SE(3)\) 几何与动量基线、端口哈密顿风格基线参照 \(\SE(3)\) 哈密顿及李群端口哈密顿神经常微分方程的思路~\citep{duong2021se3hamnode,duong2024liegroupphnode}，但均在统一的 AUV 状态契约、数据协议与评估方式下重新实现为受控候选，而非复刻文献中的原始模型。为使对照反映结构差异而非容量差异，各候选采用容量匹配的网络宽度，海流任务统一以 REMUS~100 的质量矩阵与执行器先验初始化；族内差异因而可在容量可比的前提下归因于结构配置。表~\ref{tab:model-comparison-roles} 汇总各类别的检验目标。

\begin{table}[htbp]
  \centering
  \caption{对比模型与结构消融的检验目标}
  \label{tab:model-comparison-roles}
  \begin{tabularx}{\textwidth}{>{\raggedright\arraybackslash}p{0.26\textwidth}>{\raggedright\arraybackslash}p{0.31\textwidth}>{\raggedright\arraybackslash}X}
    \toprule
    模型类别 & 代表模型 & 检验目标 \\
    \midrule
    完整结构化模型 & AUVHamNODE & 同时包含速度契约、\(\SE(3)\) 运动学、机械存储、\(D_\theta\)/\(J_\theta\)/\(\tau_\theta\) 分解和执行器滞后，作为结构化假设空间的完整候选项。 \\
    端口哈密顿风格基线 & 合并非保守广义力基线、位形广义力基线 & 保留机械存储核心时，比较不同非保守广义力组织方式对长期递推的影响。 \\
    \(\SE(3)\) 加速度黑箱 & 显式 \(\SE(3)\) 运动学与黑箱加速度动力学 & 保留位姿几何但不保留完整机械存储分解时的对照。 \\
    \(\SE(3)\) 动量黑箱 & 显式 \(\SE(3)\) 运动学、常质量映射与黑箱动量动力学 & 保留动量坐标与质量映射、动量动力学交由黑箱时的对照。 \\
    全状态黑箱 & 非结构化状态导数模型 & 缺少显式位姿几何与机械能量结构时的长期递推对照。 \\
    结构消融 & 无质量先验、对角耗散、无零功率升力耦合、广义力仅执行器条件化 & 对质量先验、耦合耗散、斜对称零功率耦合与广义力条件化范围的逐项对照。 \\
    \bottomrule
  \end{tabularx}
\end{table}

结构消融在完整模型上逐项移除或弱化单个结构组件，每个消融对应一条可检验的结构假设，并与第~\ref{sec:energy-power-properties} 节的相应性质相对接。质量先验消融移除基于物理质量矩阵的先验初始化（同时保留正定质量参数化），检验质量先验对动量--速度配对学习与长期递推的作用；对角阻尼消融把半正定耦合耗散矩阵 \(D_\theta(\xi)\) 限制为非负对角形式，检验跨通道耗散耦合相对独立通道耗散的作用；零功率耦合消融移除斜对称分支 \(J_\theta(\xi)\nu_r\)，检验升力或横向水动力类零功率耦合在候选模型族内的作用，但不等同于断言真实水动力升力项被唯一辨识；广义力条件化消融把广义力分支 \(\tau_\theta\) 收缩为主要由执行器状态或命令驱动、弱化相对速度与海流条件化，检验开放广义力通道条件化范围的作用。

所有候选在统一的数据、训练与数值实现配置下比较\footnote{连续时间积分基于 torchdiffeq~\citep{chen2018neuralode} 的固定步长四阶 Runge--Kutta 实现。}，主要设置汇总于表~\ref{tab:experimental-config}，命令级参数与运行元数据保留在各运行记录中。进入定量比较的运行须共享同一数据协议、代码环境与评估方式；出现训练崩溃或长期递推发散的运行移出精度聚合并单独披露，其中全部种子均发散者不报告中位数、而记为长期稳定性失败，以与精度不佳相区别。\(N\) 为排除后参与聚合的有效种子数；异常的判定与披露规则见第~\ref{sec:results-structure-evidence} 节。

\begin{table}[htbp]
  \centering
  \caption{长期递推实验的主要配置}
  \label{tab:experimental-config}
  \begin{tabularx}{\textwidth}{>{\raggedright\arraybackslash}p{0.18\textwidth}>{\raggedright\arraybackslash}X}
    \toprule
    配置维度 & 设置 \\
    \midrule
    数据规模 & 海流数据集含一千条轨迹、每条一百五十个控制块；控制块时长 0.2~s，状态采样间隔 0.05~s，训练与验证按 \(0.8\) 划分。 \\
    速度契约 & 数据态保存总体速度，训练前按式~\eqref{eq:velocity-contract} 转换为模型态相对水速度。 \\
    数值积分 & 固定步长四阶 Runge--Kutta，单精度浮点状态；训练、块级评估与长期递推共用同一求解器。 \\
    网络容量 & 各候选采用容量匹配的网络宽度；海流任务以 REMUS 质量矩阵与执行器先验初始化。 \\
    优化 & AdamW，学习率经线性预热后余弦衰减，至多三百个 epoch；批量、学习率端点与正则权重记录于运行配置。 \\
    递推评估 & 重采样初值，覆盖 PRBS、CHIRP 与 OU 三类控制场景，每场景 30 条评估轨迹；评估时域取 10、30 与 60~s，主时域 60~s。 \\
    \bottomrule
  \end{tabularx}
\end{table}

\subsection{初值扰动与鲁棒性评估}
\label{subsec:ic-perturbation}

初值扰动用于检验模型对带噪初始状态的鲁棒性，分别在训练与评估两个阶段引入：训练阶段作为鲁棒性正则，评估阶段作为鲁棒性测试。两阶段是不同的干预——前者决定模型是否在带噪初值上学习，后者决定模型在何种扰动下受考核——在分析中不应混为一谈。

噪声只作用于 \(t=0\) 的初始状态，并在常微分方程语义上保持一致：由干净数据态转为干净模型态后，在相对水速度 \(\nu_r\)、姿态 \(R\)、执行器反馈 \(u_a\) 与海流 \(v_c\) 通道上采样零均值扰动，再启动递推。如此控制的是模型实际消费的变量误差预算，并在海流设置下依式~\eqref{eq:velocity-contract} 保持 \(R\)、\(\nu_r\) 与 \(v_c\) 的语义一致，避免数据态扰动看似不大、而模型态输入已被过度污染的失配。该接口的定位是滤波后状态的鲁棒性正则，而非完整导航误差后验：各通道误差预算主要由 profile 常数与训练集统计量确定，目标是约束模型消费变量的误差半径，而非复现导航滤波器的联合协方差结构。

各通道预算以 REMUS~100 的航位推算与惯性量级为参照，列于表~\ref{tab:noise-budget}；评估按强度递增分为名义与退化两档，并设一档在名义随机扰动上叠加固定航向偏置的有偏配置。相对速度噪声按通道自然尺度取 \(\sigma_i=\max(\text{floor}_i,\ \alpha\,\mathrm{std}_i)\)，其中 \(\mathrm{std}_i\) 取自训练集统计、\(\alpha\) 由 profile 决定；由此，名义与退化配置的部分语义是相对训练数据分布定义的，而非完全固定的绝对物理规格——同一数据集内的跨模型比较因而合理，跨数据集时实际强度则可能漂移，各运行遂记录最终落地的预算以备追溯。除扰动强度外，噪声在采样方式上还分为逐控制块独立与同一轨迹共享一次实现两种结构，二者在训练与评估两端均可实例化；其作为训练噪声变体对长期递推的影响，在第~\ref{sec:results-structure-evidence} 节考察。

\begin{table}[htbp]
  \centering
  \caption{初值扰动各通道的 profile 噪声预算（REMUS~100 航位推算/惯性量级参照）}
  \label{tab:noise-budget}
  \small
  \begin{tabularx}{\textwidth}{>{\raggedright\arraybackslash}X cccc}
    \toprule
    通道（参数） & 名义训练 & 名义评估 & 退化评估 & 航向偏置评估 \\
    \midrule
    相对线速度 \(\nu_r\)（\(\alpha\)）& 0.03 & 0.05 & 0.10 & 同名义 \\
    相对角速度 \(\nu_r\)（\(\alpha\)）& 0.03 & 0.05 & 0.10 & 同名义 \\
    姿态（yaw 主导小角 / rad）& 0.0035 & 0.0050 & 0.0120 & 名义 \(+\) 固定 0.015 偏置 \\
    海流 \(v_{cx},v_{cy}\) / \(v_{cz}\)（m/s）& 0.008 / 0.004 & 0.012 / 0.006 & 0.030 / 0.015 & 同名义 \\
    执行器 \(\delta_r,\delta_s\)（rad）/ 转速（rpm）& 0.002 / 3 & 0.003 / 5 & 0.008 / 15 & 同名义 \\
    \bottomrule
  \end{tabularx}
  \par\vspace{2pt}
  \footnotesize 注：相对速度以 \(\sigma_i=\max(\text{floor}_i,\alpha\,\mathrm{std}_i)\) 取值，表列 \(\alpha\)；线速度与角速度下限分别为 \(0.005~\mathrm{m/s}\) 与 \(0.0015~\mathrm{rad/s}\)。航向偏置评估在名义评估的随机扰动上叠加固定 \(0.015~\mathrm{rad}\) 的 yaw 偏置，符号按样本固定、一条递推内不变。
\end{table}

训练端对噪声强度采用预热--爬坡--混合的渐进调度，以避免噪声初值在优化初期破坏训练稳定性，具体预热轮数、爬坡区间与混合比例记录于运行配置；评估端对所有模型固定噪声随机种子，使被扰动初值在不同模型之间可直接比较。
